% Imporditud: tools/import_eksperimendid.py
% Allikas: Eesti füüsikaolümpiaadi eksperimendi ülesannete
%   kogu (Rael Kalda, 2023), ülesanne Ü77, lahendus L77.
\setAuthor{EFO žürii}
\setRound{lõppvoor}
\setYear{1998}
\setNumber{G E1}
\setDifficulty{3}
\setTopic{dünaamika}
\setVahendid{niit, statiiv, kaaluviht, kell, joonlaud, millimeeterpaber graafiku jaoks}
\setPunktid{10}
\setAllikas{Eksperimendi kogu Ü77, lahendus lk 67}
\setLahendusseis{toores}

\prob{Pendel}
Valmistada niidist ja kaaluvihist pendel. Määrata pendli võnkeperioodid niidi erinevatel pikkustel. Esitada tulemus graafikuna. Leida matemaatiline seos võnkeperioodi ja pendli pikkuse vahel.

\hint

\solu
Teooria: Katse kirjeldus - mida ja kuidas mõõdab. Oluline on mitme perioodi mõõtmine ja pika niidi kasutamine, et saada suuremat täpsust (2 p).

Katse: Tulemuste esitamine tabelina (1 p), õige tulemus, st pendli pikkuse kasvades periood suureneb (1 p), graafiku telgede õige valik koos ühikutega (2 p). T ja l vahelise seose leidmise idee - graafiku lineariseerimine (1 p). Õigesti teisendatud graafik (1 p). Õige seose leidmine teisendatud graafikult, mitte valemi teadmise abil (1 p).

Järeldus: Katse tulemuse kokkuvõte, hinnang täpsusele või meetodile jne (1p).
\probend
